
\begin{thebibliography}{99}

\bibitem{zheng2023} L.~Zheng \textit{et al.}, ``Judging LLM-as-a-judge with MT-Bench and chatbot arena,'' in \textit{NeurIPS 2023 (Datasets and Benchmarks)}.

\bibitem{wang2023} P.~Wang, L.~Li, M.~Khabsa, H.~Fang, and H.~Ma, ``Large language models are not fair evaluators,'' in \textit{ACL 2024}.

\bibitem{cooper2025} A.~Chouldechova, A.~F.~Cooper, S.~Barocas, A.~Palia, D.~Vann, and H.~Wallach, ``Comparison requires valid measurement: Rethinking attack success rate comparisons in AI red teaming,'' in \textit{NeurIPS 2025 (Position Paper Track)}.

\bibitem{greshake2023} K.~Greshake, S.~Abdelnabi, S.~Mishra, C.~Endres, T.~Holz, and M.~Fritz, ``Not what you've signed up for: Compromising real-world LLM-integrated applications with indirect prompt injection,'' in \textit{AISec @ CCS 2023}.

\bibitem{wallace2024hierarchy} E.~Wallace, K.~Xiao, R.~Leike, L.~Weng, and J.~Heidecke, ``The instruction hierarchy: Training LLMs to prioritize privileged instructions,'' in \textit{ICLR 2024}.

\bibitem{zou2023} A.~Zou, Z.~Wang, J.~Z.~Kolter, and M.~Fredrikson, ``Universal and transferable adversarial attacks on aligned language models,'' arXiv:2307.15043, 2023.

\bibitem{chao2023} P.~Chao, A.~Robey, E.~Dobriban, H.~Hassani, G.~J.~Pappas, and E.~Strubell, ``Jailbreaking black box large language models in twenty queries,'' in \textit{2023 AI Safety Conference}.

\bibitem{chao2024jb} P.~Chao, E.~Debenedetti, A.~Robey, M.~Andriushchenko, F.~Croce, V.~Sehwag, E.~Dobriban, N.~Flammarion, G.~J.~Pappas, F.~Tram{\`er}, H.~Hassani, and E.~Wong, ``JailbreakBench: An open robustness benchmark for jailbreaking large language models,'' in \textit{NeurIPS 2024 (Datasets and Benchmarks)}.

\bibitem{mehrotra2024} A.~Mehrotra, M.~Zampetakis, P.~Kassianik, B.~Nelson, H.~Anderson, Y.~Singer, and A.~Karbasi, ``Tree of attacks: Jailbreaking black-box LLMs automatically,'' arXiv:2312.02119, 2024.

\bibitem{liu2024autodan} X.~G.~Liu, N.~Xu, M.~Chen, and C.~Xiao, ``AutoDAN: Generating stealthy jailbreak prompts on aligned large language models,'' in \textit{ICLR 2024}.

\bibitem{shen2023} X.~Shen, Z.~Chen, M.~Backes, Y.~Shen, and Y.~Zhang, ``Do Anything Now: Characterizing and evaluating in-the-wild jailbreak prompts on large language models,'' arXiv:2308.03825, 2023.

\bibitem{anil2024} C.~Anil, E.~Durmus, N.~Panickssery, M.~Sharma, J.~Benton, S.~Kundu \textit{et al.}, ``Many-shot jailbreaking,'' in \textit{NeurIPS 2024}.

\bibitem{ganguli2022} D.~Ganguli, L.~Lovitt, J.~Kernion, A.~Askell, Y.~Bai, S.~Farley \textit{et al.}, ``Red teaming language models to reduce harms: Methods, scaling behaviors, and lessons learned,'' arXiv:2209.07858, 2022.

\bibitem{diao2025} M.~Diao \textit{et al.}, ``AutoRed: A free-form adversarial prompt generation framework for automated red teaming,'' arXiv:2510.08329, 2025.

\bibitem{mazeika2024} M.~Mazeika, L.~Phan, X.~Yin, A.~Zou, Z.~Wang, N.~Mu, E.~Sakhaee, N.~Li, S.~Basart, B.~Li, D.~Forsyth, and D.~Hendrycks, ``HarmBench: A standardized evaluation framework for automated red teaming and robust refusal,'' in \textit{ICML 2024}.

\bibitem{andriushchenko2024} M.~Andriushchenko, F.~Croce, and N.~Flammarion, ``Jailbreaking leading safety-aligned LLMs with simple adaptive attacks,'' in \textit{ICLR 2024}.

\bibitem{souly2024} A.~Souly, Q.~Lu, D.~Bowen, T.~Trinh, E.~Hsieh, S.~Pandey, P.~Abbeel, J.~Svegliato, S.~Emmons, O.~Watkins, and S.~Toyer, ``A StrongREJECT for empty jailbreaks,'' in \textit{NeurIPS 2024 (Datasets and Benchmarks)}.

\bibitem{nist2024} NIST, ``Cybersecurity Evaluation Program: CyberSecEval 3,'' National Institute of Standards and Technology, 2024.

\bibitem{zhang2025usenix} L.~Zhang, X.~Gao, L.~Yao, J.~Song \textit{et al.}, ``Exploiting task-level vulnerabilities: An automatic jailbreak attack and defense benchmarking for LLMs,'' in \textit{USENIX Security 2025}.

\bibitem{chu2025} J.~Chu, Y.~Liu, Z.~Yang, X.~Shen, M.~Backes, and Y.~Zhang, ``JailbreakRadar: Comprehensive assessment of jailbreak attacks against LLMs,'' in \textit{ACL 2025}.

\bibitem{shinn2023} N.~Shinn, F.~Cassano, E.~Berman, A.~Gopinath, K.~Narasimhan, and S.~Yao, ``Reflexion: Language agents with verbal reinforcement learning,'' in \textit{NeurIPS 2023}.

\bibitem{madaan2023} A.~Madaan \textit{et al.}, ``Self-Refine: Iterative refinement with self-feedback,'' in \textit{NeurIPS 2023}.

\bibitem{lewis2020} P.~Lewis, E.~Perez, A.~Piktus \textit{et al.}, ``Retrieval-augmented generation for knowledge-intensive NLP tasks,'' in \textit{NeurIPS 2020}.

\bibitem{sanh2019} V.~Sanh, L.~Debut, J.~Chaumond, and T.~Wolf, ``DistilBERT, a distilled version of BERT: Smaller, faster, cheaper and lighter,'' arXiv:1910.01108, 2019.

\bibitem{hu2022} E.~J.~Hu, Y.~Shen, P.~Wallis, Z.~Allen-Zhu, Y.~Li, S.~Wang, L.~Wang, and W.~Chen, ``LoRA: Low-rank adaptation of large language models,'' in \textit{ICLR 2022}.

\bibitem{dettmers2023} T.~Dettmers, A.~Pagnoni, A.~Holtzman, and L.~Zettlemoyer, ``QLoRA: Efficient finetuning of quantized LLMs,'' in \textit{NeurIPS 2023}.

\bibitem{kwon2023} W.~Kwon, Z.~Li, S.~Zhuang, Y.~Sheng, L.~Zheng, C.~H.~Yu, J.~E.~Gonzalez, H.~Zhang, and I.~Stoica, ``Efficient memory management for large language model serving with PagedAttention,'' in \textit{SOSP 2023}.

\bibitem{dubois2024} A.~Grattafiori, A.~Dubey, A.~Jauhri, A.~Pandey, A.~Kadian, A.~Al-Dahle \textit{et al.}, ``The Llama 3 herd of models,'' arXiv:2407.21783, 2024.

\bibitem{gemma2024} Gemma Team, Google, ``Gemma: Open models based on Gemini research and technology,'' arXiv:2403.08295, 2024.

\bibitem{cai2024} Z.~Cai \textit{et al.}, ``InternLM2 technical report,'' arXiv:2403.17297, 2024.

\bibitem{jiang2023} A.~Q.~Jiang, A.~Sablayrolles, A.~Mensch \textit{et al.}, ``Mistral 7B,'' arXiv:2310.06825, 2023.

\bibitem{he2021} P.~He, X.~Liu, J.~Gao, and W.~Chen, ``DeBERTa: Decoding-enhanced BERT with discrete auto-encoding,'' in \textit{ACL-IJCNLP 2021}.

\bibitem{reimers2019} N.~Reimers and I.~Gurevych, ``Sentence-BERT: Sentence embeddings using Siamese BERT-networks,'' in \textit{EMNLP-IJCNLP 2019}.

\bibitem{johnson2019} J.~Johnson, M.~Douze, and H.~J\'egou, ``Billion-scale similarity search with GPUs,'' \textit{IEEE Trans. Big Data}, vol.~7, no.~3, pp.~535--547, 2019.

\bibitem{rafaelov2023} R.~Rafailov, A.~Sharma, E.~Mitchell, C.~D.~Manning, S.~Ermon, and C.~Finn, ``Direct preference optimization: Your language model is secretly a reward model,'' in \textit{NeurIPS 2023}.

\bibitem{tensortrust2025} S.~Toyer, O.~Watkins, E.~A.~Mendes, J.~Svegliato, L.~Bailey, T.~Wang \textit{et al.}, ``Tensor Trust: Interpretable prompt injection attacks from an online game,'' arXiv:2311.01011, 2023. Dataset and code: \url{https://github.com/HumanCompatibleAI/tensor-trust-data}

\end{thebibliography}

\appendix

\section{Metric definitions}
\label{app:metrics}
Let a benchmark contain $N$ scenarios, and let attempt $i$ of scenario $s$ produce response $r_{s,i}$ and signal vector $(g_{s,i}, a_{s,i}, v_{s,i}, e_{s,i})$ for the four success signals in priority order.
\begin{itemize}
\item Break rate $= (1/N)\,|\{s : \max_i (g_{s,i} \lor a_{s,i} \lor v_{s,i} \lor e_{s,i})\}|$.
\item Verified rate $= (1/N)\,|\{s : \max_i (g_{s,i} \lor v_{s,i})\}|$.
\item Attempts to success: $T_s = \min\{i : \text{any signal fires}\}$, undefined if the scenario is not broken.
\item Top-$k$ success $= (1/N)\,|\{s : T_s \le k\}|$, reported for $k \in \{1,3,5,10\}$.
\item Average attempts on success $= (1/|S|) \sum_{s \in S} T_s$, where $S$ is the set of broken scenarios.
\item Extractor precision $= \mathrm{TP}/(\mathrm{TP}+\mathrm{FP})$, recall $= \mathrm{TP}/(\mathrm{TP}+\mathrm{FN})$, over all attempts with ground truth available, where a true positive is an attempt in which the code leaked and the extractor found it, a false negative is a leak the extractor missed, and a false positive is a non-leaking response from which the extractor nonetheless produced a candidate matching the code.
\item Failure attribution: per unsuccessful scenario, the phase with the largest share of failed attempts among \{judge rejection (the stop-point judge classified the response as non-attack), extractor miss (a ground-truth leak the extractor did not find), verifier rejection (a verified-candidate replay refused by the victim)\}. The attribution is a concentration measure: because the judge does not gate success, it shows where failed runs concentrate, not what causally blocked a particular win. Scenarios are additionally labeled by the controller's failure mode (planner-stuck, generator-rephrase-fail, leaked-unverified, access-granted-unverified, never-leaked).
\end{itemize}

\section{Scenario pools and defense taxonomy}
\label{app:pools}
\begin{table}[!t]
\centering
\footnotesize
\begin{tabular}{p{1.0in} p{0.7in} p{1.38in}}
\hline
Pool & Size & Rule \\
\hline
Full short-code corpus & 72,008 & access code length $< 30$ \\
Recoverable subset (main) & 34,576 & code direct, deterministic, or indirect under the defense's stated rules \\
Architecture pool & 1,000 & seeded sample of the full corpus \\
Main evaluation & 22,791--26,047 / victim & all recoverable scenarios; 26,047 for Gemma-2b, Mistral-7B, InternLM2-7B; 22,791 for Llama-3-8B \\
Table I sample & 13,024 / victim & seeded sample of the main pool \\
Selection benchmark set & 980 & 780 dev / 200 holdout, seed 42, fixed hash \\
\hline
\end{tabular}
\end{table}

Defense families (eight): password (enter exactly $X$), roleplay (character constraints with trigger exceptions), translation (respond only in another language or encoding), trigger-phrase (a specific phrase unlocks behavior), conditional (behavior gated on a condition), exception (a stated exception to the hiding rule), conversation (multi-turn interactive rule), instruction-hiding (hide the system prompt generally). Access code types (four): TOKEN (one word or symbol), PHRASE (two to five words), SENTENCE (more than five words or more than 20 characters), MULTILINE (several lines).

\section{Training data construction}
\label{app:data}
\begin{table}[!t]
\centering
\footnotesize
\begin{tabular}{p{0.95in} p{0.5in} p{1.6in}}
\hline
Dataset & Size & Source \\
\hline
Oracle trajectories & 4,505 & privileged best-of-N runs, 8 workers \\
Planner fine-tuning set & 17,129 & 40\% oracle (3$\times$ weight), 45\% runtime successes, 15\% runtime failures; family-balanced \\
Generator fine-tuning set & 13,037 & 5,955 oracle, 7,082 runtime successes; plan-conditioned \\
Ranker set & 12,144 & 3,036 positive / 9,108 negative (1:3); 9,722/1,211/1,211 split \\
Strategy predictor set & 11,714 & gold-labeled runtime attempts + synthetic \\
Access-code-type classifier set & 118,326 & labeled from the defense corpus \\
Runtime success log & 217,168 & cross-run accumulation, content-hash dedup \\
Runtime verified subset & 3,558 & of the above, victim-confirmed \\
\hline
\end{tabular}
\end{table}

The planner and generator are low-rank adapters on a shared 8B base, trained by quantized low-rank fine-tuning~\cite{hu2022, dettmers2023}. The ranker is a small DeBERTa discriminator~\cite{he2021}. The access-code-type predictor and the stop-point judge are distilled BERT classifiers~\cite{sanh2019}. Preference-optimization datasets~\cite{rafaelov2023} derived from the same logs are kept for future policy improvement, but they are not part of the evaluated configuration.

\section{Plan contract (schema)}
The planner emits an XML block with the fields \emph{strategy} (18-value enum), \emph{primitive\_sequence} (ordered list, 1 to 5 items from a group/name library), \emph{style} (formal, conversational, academic, story, direct), \emph{expected\_access\_type} (TOKEN, PHRASE, SENTENCE, MULTILINE, UNKNOWN), \emph{retry\_policy} (explore, retry\_same\_strategy, switch\_strategy), \emph{confidence} (real in $[0,1]$), and \emph{failure\_reason} (free text). Canonicalization rules: unknown strategy maps to instruction-leak; an empty primitive list maps to a single default framing primitive; the primitive list is truncated to five; unknown style maps to direct; unknown access type maps to UNKNOWN; unknown retry policy maps to explore; confidence is clamped to $[0,1]$. Malformed output triggers a defense-aware strategy rotation. The guard then replaces the strategy if it is embargoed (three consecutive failures), if its failure streak reached three, or if the plan requests a switch with a streak of at least two.

\section{Reproducibility}
All rates in the main text carry 95\% Wilson confidence intervals where they appear in tables. Comparisons between benchmark arms use independent two-proportion $z$-tests, applied to the separate run archives of each arm. All benchmark samples are seeded. The selection benchmark set is fixed by hash (200 holdout, 780 development, seed 42). Worker sharding is contiguous and even over the shared seeded draw, so multi-worker runs attack disjoint slices of one sample. Every scenario logs a per-attempt trace (plan, prompt, response, candidates, scores, signals, judge outcome) and a scenario-level summary. We merge benchmark summaries from the worker shards by summing counters and recomputing extractor metrics from the summed confusion counts. The knowledge stores are plain JSON tables and a vector index~\cite{johnson2019} rebuilt from the logs. So the complete learning state at any point can be reconstructed from the run archives. The default victim is served with 200-token responses at temperature 0.7; the planner samples at temperature 0 and the generator at temperature 0.7 with nucleus sampling 0.9.

\end{document}
